\documentclass[12pt]{article}

% --------------------
% Packages
% --------------------
\usepackage[margin=1in]{geometry}
\usepackage{amsmath, amssymb}
\usepackage{booktabs}
\usepackage{graphicx}
\usepackage{setspace}
\usepackage{hyperref}
\usepackage{listings}
\usepackage{xcolor}

% --------------------
% Formatting
% --------------------
\onehalfspacing

\hypersetup{
    colorlinks=true,
    linkcolor=blue,
    urlcolor=blue,
    citecolor=blue
}

% Code style
\lstset{
    basicstyle=\ttfamily\small,
    breaklines=true,
    frame=single,
    columns=fullflexible
}

% --------------------
% Title
% --------------------
\title{Running Simple Regressions with Penn World Table Data in \texttt{R}}
\author{Harlly Zhou}
\date{\today}

% --------------------
% Document
% --------------------
\begin{document}

\maketitle
\footnote{This tutorial note is written with the help of AI.}

\subsection*{1. Basic \texttt{R} operations}

Before working with data, you should understand a few basic commands.

\begin{itemize}
    \item Check current working directory:
\begin{verbatim}
getwd()
\end{verbatim}

    \item Set working directory (modify the path to your folder):
\begin{verbatim}
setwd("YOUR/PATH/HERE")
\end{verbatim}

    \item List files in the folder:
\begin{verbatim}
list.files()
\end{verbatim}

    \item Install packages (run once):
\begin{verbatim}
install.packages("dplyr")
install.packages("readr")
install.packages("car")
\end{verbatim}

    \item Load packages (run every time):
\begin{verbatim}
library(dplyr)
library(readr)
library(car)
\end{verbatim}
\end{itemize}

\textbf{Important:} Place your \texttt{pwt.csv} file in the working directory.

---

\subsection*{2. Load and inspect the data}

\begin{verbatim}
pwt <- read_csv("pwt.csv")

# Check column names
names(pwt)

# Preview data
head(pwt)

# Check available countries
unique(pwt$country)
\end{verbatim}

---

\subsection*{3. Select countries and years}

Fill in the countries (already given in the question):

\begin{verbatim}
countries <- c("United States", "Japan", "India", "Nigeria")

pwt_sub <- pwt %>%
  filter(country %in% countries,
         year >= 1990,
         year <= 2020)
\end{verbatim}

---

\subsection*{4. Construct variables}

Fill in the correct PWT variable names (from the question):

\begin{verbatim}
pwt_sub <- pwt_sub %>%
  transmute(
    country = country,
    year = year,
    Y = ______ * 1e6,   # output
    K = ______ * 1e6,   # capital
    N = ______ * 1e6    # labor
  ) %>%
  filter(Y > 0, K > 0, N > 0) %>%
  arrange(country, year) %>%
  group_by(country) %>%
  mutate(
    lnY = log(Y),
    lnK = log(K),
    lnN = log(N),
    gY  = lnY - lag(lnY),
    gK  = lnK - lag(lnK),
    gN  = lnN - lag(lnN)
  ) %>%
  ungroup()
\end{verbatim}

\textbf{Check:} Why do we use \texttt{group\_by(country)} here?

---

\subsection*{5. Level regression (Task 1)}

\paragraph{Step 1: Run regression for one country}

Fill in the country name:

\begin{verbatim}
df <- pwt_sub %>%
  filter(country == "________")

reg_level <- lm(lnY ~ lnK + lnN, data = df)

summary(reg_level)
\end{verbatim}

\paragraph{Step 2: Extract coefficients}

\begin{verbatim}
coef(reg_level)
\end{verbatim}

\textbf{Task:} Identify $\hat{\beta}_1$ and $\hat{\beta}_2$.

---

\paragraph{Step 3: Sum of coefficients}

\begin{verbatim}
b <- coef(reg_level)

b["lnK"] + b["lnN"]
\end{verbatim}

\textbf{Task:} Is this close to 1?

---

\paragraph{Step 4: $F$-test}

\begin{verbatim}
linearHypothesis(reg_level, "lnK + lnN = 1")
\end{verbatim}

\textbf{Task:} Look at the $p$-value. Do you reject $H_0$?

---

\paragraph{Step 5: Repeat for all countries}

Replace the country name each time:
\begin{verbatim}
"United States"
"Japan"
"India"
"Nigeria"
\end{verbatim}

---

\subsection*{6. Growth regression (Task 2)}

\paragraph{Step 1: Prepare growth data}

\begin{verbatim}
df_g <- df %>%
  filter(!is.na(gY), !is.na(gK), !is.na(gN))
\end{verbatim}

Why do we drop missing values here?

---

\paragraph{Step 2: Run regression}

\begin{verbatim}
reg_growth <- lm(gY ~ gK + gN, data = df_g)

summary(reg_growth)
\end{verbatim}

---

\paragraph{Step 3: Extract coefficients}

\begin{verbatim}
coef(reg_growth)
\end{verbatim}

\textbf{Task:} Identify $\hat{\gamma}_1$ and $\hat{\gamma}_2$.

---

\paragraph{Step 4: Sum of coefficients}

\begin{verbatim}
g <- coef(reg_growth)

g["gK"] + g["gN"]
\end{verbatim}

---

\paragraph{Step 5: $F$-test}

\begin{verbatim}
linearHypothesis(reg_growth, "gK + gN = 1")
\end{verbatim}

---

\subsection*{7. Loop version (optional)}

To avoid repeating code, you may use a loop:

\begin{verbatim}
for (cty in countries) {

  df <- pwt_sub %>% filter(country == cty)

  reg_level <- lm(lnY ~ lnK + lnN, data = df)

  df_g <- df %>% filter(!is.na(gY))

  reg_growth <- lm(gY ~ gK + gN, data = df_g)

  print(cty)
  print(coef(reg_level))
  print(coef(reg_growth))
}
\end{verbatim}

---

\subsection*{8. Interpretation hints}

When answering the questions:

\begin{itemize}
    \item Coefficients may not lie between 0 and 1.
    \item The sum may differ from 1.
    \item $F$-tests may reject the theoretical restriction.
\end{itemize}

These results suggest that the regression does not perfectly recover the theoretical parameter $\alpha$.

Think about:
\begin{itemize}
    \item omitted variables,
    \item endogeneity,
    \item data limitations.
\end{itemize}

---

\subsection*{9. Summary}

This tutorial shows how to:
\begin{itemize}
    \item manage files and directories in \texttt{R},
    \item load and inspect data,
    \item construct variables,
    \item run regressions,
    \item test hypotheses.
\end{itemize}

More importantly, it shows that theoretical parameters are not always directly estimated from simple regressions.


\end{document}